% Part: first-order-logic
% Chapter: syntax-and-semantics
% Section: satisfaction

\documentclass[../../../include/open-logic-section]{subfiles}

\begin{document}

\olfileid[fr]{fol}{syn}{sat}

\olsection{Satisfaction des \printtoken{p}{formula}
  dans les \printtoken{p}{structure}}

\begin{explain}
Les notions fondamentales qui relient les expressions, notamment les
termes et les !!{formula}s, aux !!{structure}s sont la
\emph{!!{value}} d'un terme et la \emph{satisfaction} d'!!a{formula}.
Informellement, la !!{value} d'un terme est un !!{element} d'!!a{structure} :
si le terme est simplement une constante, sa !!{value} est l'objet
que la !!{structure} attribue à cette constante; s'il est formé à
l'aide de !!{function}s, sa !!{value} se calcule à partir des
!!{value}s des constantes et des fonctions que la !!{structure}
attribue aux symboles figurant dans le terme. Une !!{formula} est
\emph{satisfaite} dans !!a{structure} si l'interprétation des
!!{predicate}s rend la !!{formula} vraie dans le !!{domain} de la
!!{structure}. La satisfaction est définie inductivement : la
!!{structure} détermine directement quand les !!{formula}s atomiques
sont satisfaites; pour une !!{formula} complexe, la définition dépend
de son connecteur ou quantificateur principal et de la satisfaction,
ou non, de ses !!{subformula}s immédiates.

Le cas des quantificateurs est plus délicat : la !!{subformula}
immédiate d'une !!{formula} quantifiée peut contenir une !!{variable}
libre, dont la !!{value} n'est pas spécifiée par la !!{structure}.\footnote{Précision éditoriale : la source affirme la présence d'une variable libre. La quantification peut cependant être vacuante, puisque les règles de formation n'exigent pas que la variable quantifiée apparaisse dans la formule.}
Pour traiter cette difficulté, nous introduisons des
\emph{assignations de variables} et définissons la satisfaction
relativement à la fois à !!a{structure} et à une assignation de
!!{variable}s.
\end{explain}

\begin{defn}[Assignation de variables]
Une \emph{assignation de variables}~$s$ pour !!a{structure}~$\Struct{M}$
est une fonction qui associe à chaque !!{variable} un !!{element}
de~$\Domain M$, autrement dit $s\colon \Var \to \Domain M$.
\end{defn}

\begin{explain}
Une !!{structure} attribue une !!{value} à chaque !!{constant}, et une
assignation attribue une valeur à chaque !!{variable}. Nous voulons aussi que
les termes construits à partir de ces symboles désignent des
!!{element}s du !!{domain}. Nous définissons donc inductivement la
!!{value} des termes. Pour une !!{constant} ou une variable, elle
est celle que spécifie respectivement la !!{structure} ou l'assignation;
pour un terme plus complexe, elle se calcule récursivement au moyen des
fonctions attribuées aux !!{function}s par la !!{structure}.
\end{explain}

\begin{defn}[!!^{value} des termes]
Si $t$ est un terme du langage~$\Lang L$, si $\Struct M$ est une
!!{structure} pour~$\Lang L$ et si $s$ est une assignation de
!!{variable}s pour~$\Struct M$, la \emph{!!{value}}~$\Value{t}{M}[s]$
est définie comme suit :
\begin{enumerate}
\item \indcase{t}{c}{$\Value{\indfrm}{M}[s] = \Assign{\indcomplex}{M}$.}
\item \indcase{t}{x}{$\Value{\indfrm}{M}[s] = s(\indcomplex)$.}
\item \indcase{t}{\Atom{f}{t_1, \ldots, t_n}}{
\[
\Value{\indfrm}{M}[s] = \Assign{f}{M}(\Value{t_1}{M}[s], \ldots,
\Value{t_n}{M}[s]).
\]}
\end{enumerate}
\end{defn}

\begin{defn}[Variante en~$x$]
Soit $s$ une assignation de !!{variable}s pour !!a{structure}~$\Struct M$.
Toute assignation~$s'$ pour~$\Struct M$ qui ne diffère de~$s$ que,
éventuellement, par la !!{value} attribuée à~$x$ est appelée
\emph{variante en~$x$} de~$s$. Si $s'$ est une telle variante,
nous écrivons $\varAssign{s'}{s}{x}$.
\end{defn}

\begin{explain}
Une variante en~$x$ de l'assignation~$s$ n'attribue pas
\emph{nécessairement} une valeur différente à~$x$. Toute assignation
est donc une variante en~$x$ d'elle-même.
\end{explain}

\begin{defn}
  Si $s$ est une assignation de !!{variable}s pour
  !!a{structure}~$\Struct M$, si $x$ est une !!{variable} et si
  $m \in \Domain{M}$, alors l'assignation~$\Subst{s}{m}{x}$ est définie par
  \[\Subst{s}{m}{x}(y) = \begin{cases}
    m & \text{si } y \ident x\\
    s(y) & \text{sinon}.
  \end{cases}\]
\end{defn}

Autrement dit, $\Subst{s}{m}{x}$ est la variante en~$x$ de~$s$ qui
attribue à~$x$ l'!!{element}~$m$ du !!{domain} et qui conserve les
valeurs attribuées par~$s$ aux autres !!{variable}s.

\begin{defn}[Satisfaction]
\ollabel{defn:satisfaction}
La satisfaction d'!!a{formula}~$!A$ dans !!a{structure}~$\Struct M$
relativement à une assignation de !!{variable}s~$s$, notée
$\Sat{M}{!A}[s]$, est définie récursivement comme suit. (Nous
écrivons $\Sat/{M}{!A}[s]$ pour signifier « $\Sat{M}{!A}[s]$
ne vaut pas ».)
\begin{enumerate}
\tagitem{prvFalse}{%
  \indcase{!A}{\lfalse}{$\Sat/{M}{\indfrm}[s]$.}}{}

\tagitem{prvTrue}{%
  \indcase{!A}{\ltrue}{$\Sat{M}{\indfrm}[s]$.}}{}

\item \indcase{!A}{\Atom{R}{t_1, \dots, t_n}}{$\Sat{M}{\indfrm}[s]$
  si et seulement si $\langle \Value{t_1}{M}[s], \dots, \Value{t_n}{M}[s] \rangle \in
  \Assign{R}{M}$.}

\item \indcase{!A}{\eq[t_1][t_2]}{$\Sat{M}{\indfrm}[s]$ si et seulement si
  $\Value{t_1}{M}[s] = \Value{t_2}{M}[s]$.}

\tagitem{prvNot}{%
  \indcase{!A}{\lnot !B}{$\Sat{M}{\indfrm}[s]$ si et seulement si
    $\Sat/{M}{!B}[s]$.}}{}

\tagitem{prvAnd}{%
  \indcase{!A}{(!B \land !C)}{$\Sat{M}{\indfrm}[s]$ si et seulement si
    $\Sat{M}{!B}[s]$ et $\Sat{M}{!C}[s]$.}}{}

\tagitem{prvOr}{%
  \indcase{!A}{(!B \lor !C)}{$\Sat{M}{\indfrm}[s]$ si et seulement si
    $\Sat{M}{!B}[s]$ ou $\Sat{M}{!C}[s]$ (ou les deux).}}{}

\tagitem{prvIf}{%
  \indcase{!A}{(!B \lif !C)}{$\Sat{M}{\indfrm}[s]$ si et seulement si
    $\Sat/{M}{!B}[s]$ ou $\Sat{M}{!C}[s]$ (ou les deux).}}{}

\tagitem{prvIff}{%
  \indcase{!A}{(!B \liff !C)}{$\Sat{M}{\indfrm}[s]$ si et seulement si
    $\Sat{M}{!B}[s]$ et $\Sat{M}{!C}[s]$ valent tous deux, ou si ni
    $\Sat{M}{!B}[s]$ ni $\Sat{M}{!C}[s]$ ne vaut.}}{}

\tagitem{prvAll}{%
  \indcase{!A}{\lforall[x][!B]}{$\Sat{M}{\indfrm}[s]$ si et seulement si,
    pour tout
    !!{element}~$m \in \Domain M$, $\Sat{M}{!B}[\Subst{s}{m}{x}]$.}}{}

\tagitem{prvEx}{%
  \indcase{!A}{\lexists[x][!B]}{$\Sat{M}{\indfrm}[s]$ si et seulement si,
  pour au moins un !!{element}~$m \in \Domain M$,
  $\Sat{M}{!B}[\Subst{s}{m}{x}]$.}}{}
\end{enumerate}
\end{defn}

\begin{explain}
Les assignations de variables interviennent dans
\iftag{notprvEx,notprvAll}{la dernière clause}{les deux dernières clauses}.
\iftag{prvAll}{ Nous ne pouvons pas définir la satisfaction de
$\lforall[x][!B(x)]$ en disant « pour tout $m \in \Domain{M}$,
$\Sat{M}{!B(m)}$ ».}{}\iftag{prvEx}{ Nous ne pouvons pas définir la
satisfaction de $\lexists[x][!B(x)]$ en disant « pour au moins un
$m \in \Domain{M}$, $\Sat{M}{!B(m)}$ ».}{} En effet, si
$m \in \Domain M$, alors $m$ n'est pas un symbole du langage et
$!B(m)$ n'est donc pas !!a{formula} (autrement dit,
$\Subst{!B}{m}{x}$ n'est pas défini). Nous ne pouvons pas non plus
supposer qu'il existe assez de !!{constant}s ou de termes pour
désigner tous les !!{element}s de~$\Struct{M}$ : la définition d'une
!!{structure} ne l'exige pas. Dans le langage standard, l'ensemble
des !!{constant}s est !!{denumerable}; si $\Domain{M}$ n'est pas
!!{enumerable}, il n'y a même pas assez de !!{constant}s pour nommer
chaque objet.

Nous résolvons cette difficulté en introduisant les assignations de
!!{variable}s, qui relient directement les variables aux !!{element}s
du !!{domain}. Au lieu de dire, par exemple, que
\iftag{prvEx}{$\lexists[x][!B(x)]$}{$\lforall[x][!B(x)]$} est satisfaite
dans~$\Struct M$ si et seulement si
\iftag{prvEx}{pour au moins un $m \in \Domain{M}$, $\Sat{M}{!B(m)}$}
{pour tout $m \in \Domain{M}$, $\Sat{M}{!B(m)}$}, nous disons qu'elle
est satisfaite dans~$\Struct M$ \emph{relativement à}~$s$ si et
seulement si $!B(x)$ est satisfaite relativement
à~$\Subst{s}{m}{x}$ \iftag{prvEx}{pour au moins un}{pour tout}
$m \in \Domain M$.
\end{explain}

\begin{ex}
Soit $\Lang{L} = \{a, b, f, R\}$, où $a$ et $b$ sont des
!!{constant}s, $f$ un !!{function} binaire et $R$ un !!{predicate}
binaire. Considérons la !!{structure}~$\Struct{M}$ définie par :
\begin{enumerate}
\item $\Domain M = \{1, 2, 3, 4\}$
\item $\Assign{a}{M} = 1$
\item $\Assign{b}{M} = 2$
\item $\Assign{f}{M}(x, y) = x+y$ si $x+y \le 3$, et vaut~$3$ sinon.
\item $\Assign{R}{M} = \{\tuple{1, 1}, \tuple{1, 2}, \tuple{2, 3}, \tuple{2, 4}\}$
\end{enumerate}
La fonction $s(x) = 1$, qui attribue $1 \in \Domain{M}$ à chaque
!!{variable}, est une assignation de variables pour~$\Struct{M}$.

On a alors
\begin{align*}
\Value{f(a,b)}{M}[s] & = \Assign{f}{M}(\Value{a}{M}[s], \Value{b}{M}[s]).
\intertext{Puisque $a$ et $b$ sont des !!{constant}s, $\Value{a}{M}[s]
  = \Assign{a}{M} = 1$ et $\Value{b}{M}[s] = \Assign{b}{M} = 2$. Donc}
\Value{f(a,b)}{M}[s] & = \Assign{f}{M}(1, 2) = 1+2 = 3.
\intertext{Pour calculer la valeur de $f(f(a,b),a)$, nous considérons}
\Value{f(f(a,b),a)}{M}[s] & = \Assign{f}{M}(\Value{f(a, b)}{M}[s],
\Value{a}{M}[s]) = \Assign{f}{M}(3, 1) = 3,
\intertext{car $3+1 > 3$. Comme $s(x) = 1$ et $\Value{x}{M}[s] =
  s(x)$, nous avons aussi}
\Value{f(f(a,b),x)}{M}[s] & = \Assign{f}{M}(\Value{f(a, b)}{M}[s],
\Value{x}{M}[s]) = \Assign{f}{M}(3, 1) = 3,
\end{align*}

Une !!{formula} atomique~$R(t_1, t_2)$ est satisfaite si le couple
des valeurs de ses arguments, soit $\tuple{\Value{t_1}{M}[s],
  \Value{t_2}{M}[s]}$, appartient à~$\Assign{R}{M}$. Ainsi,
$\Sat{M}{R(b,f(a,b))}[s]$, car $\tuple{\Value{b}{M},
  \Value{f(a,b)}{M}} = \tuple{2, 3} \in \Assign{R}{M}$, mais
$\Sat/{M}{R(x, f(a,b))}[s]$, car
$\tuple{1, 3} \notin \Assign{R}{M}$.

Pour déterminer si une !!{formula} non atomique~$!A$ est satisfaite,
on applique la clause de la définition inductive correspondant à son
connecteur principal. Par exemple, le connecteur principal de
$R(a, a) \lif (R(b, x) \lor R(x, b))$ est~$\lif$, d'où
\begin{align*}
  & \Sat{M}{R(a, a) \lif (R(b, x) \lor R(x, b))}[s]
    \text{ si et seulement si }\\
  & \qquad
  \Sat/{M}{R(a,a)}[s] \text{ ou } \Sat{M}{R(b, x) \lor R(x, b)}[s]
  \intertext{Comme $\Sat{M}{R(a,a)}[s]$ (car $\tuple{1,1} \in
    \Assign{R}{M}$), nous devons encore déterminer si
    $\Sat{M}{R(b, x) \lor R(x, b)}[s]$ :}
  & \Sat{M}{R(b, x) \lor R(x, b)}[s]
    \text{ si et seulement si }\\
  & \qquad \Sat{M}{R(b,x)}[s] \text{ ou } \Sat{M}{R(x, b)}[s]
  \intertext{C'est bien le cas, puisque $\Sat{M}{R(x, b)}[s]$
    (car $\tuple{1,2} \in \Assign{R}{M}$).}
\end{align*}

Rappelons qu'une variante en~$x$ de~$s$ est une assignation de
variables qui ne diffère de~$s$ que, éventuellement, par la valeur
attribuée à~$x$. Chaque !!{element} de~$\Domain{M}$ détermine une
telle variante :
\begin{align*}
  s_1 & = \Subst{s}{1}{x}, &
  s_2 & = \Subst{s}{2}{x},\\
  s_3 & = \Subst{s}{3}{x}, &
  s_4 & = \Subst{s}{4}{x}.
\end{align*}
Par exemple, $s_2(x) = 2$ et $s_2(y) = s(y) = 1$ pour toute
variable~$y$ autre que~$x$. Comme $\Domain{M} = \{1, 2, 3, 4\}$,
ce sont toutes les variantes en~$x$ de~$s$ pour~$\Struct{M}$.
En particulier, $s_1 = s$ : $s$ est toujours une variante en~$x$
d'elle-même.

\iftag{prvEx}{Pour déterminer si la !!{formula} existentielle
  $\lexists[x][!A(x)]$ est satisfaite, il faut vérifier si
  $\Sat{M}{!A(x)}[\Subst{s}{m}{x}]$ pour au moins un
  $m \in \Domain M$. Ainsi,
  \[
  \Sat{M}{\lexists[x][(R(b,x) \lor R(x,b))]}[s],
  \]
  car $\Sat{M}{R(b,x) \lor R(x, b)}[\Subst{s}{1}{x}]$
  (la variante $\Subst{s}{3}{x}$ conviendrait aussi). En revanche,
  \[
  \Sat/{M}{\lexists[x][(R(b,x) \land R(x,b))]}[s]
  \]
  car, quel que soit $m \in \Domain{M}$,
  $\Sat/{M}{R(b,x) \land R(x,b)}[\Subst{s}{m}{x}]$.}{}

\iftag{prvAll}{Pour déterminer si la !!{formula} universelle
  $\lforall[x][!A(x)]$ est satisfaite, il faut vérifier si
  $\Sat{M}{!A(x)}[\Subst{s}{m}{x}]$ pour tout $m \in \Domain M$.
  Ainsi,
  \[
  \Sat{M}{\lforall[x][(R(x,a) \lif R(a,x))]}[s],
  \]
  car $\Sat{M}{R(x,a) \lif R(a,x)}[\Subst{s}{m}{x}]$ pour tout
  $m \in \Domain M$. Si $m = 1$, alors
  $\Sat{M}{R(a,x)}[\Subst{s}{1}{x}]$ : le conséquent est vrai.
  Si $m = 2$, $3$ ou~$4$, alors
  $\Sat/{M}{R(x,a)}[\Subst{s}{m}{x}]$ : l'antécédent est faux.
  En revanche,
  \[
  \Sat/{M}{\lforall[x][(R(a,x) \lif R(x,a))]}[s]
  \]
  car $\Sat/{M}{R(a,x) \lif R(x,a)}[\Subst{s}{2}{x}]$ :
  $\Sat{M}{R(a, x)}[\Subst{s}{2}{x}]$, tandis que
  $\Sat/{M}{R(x, a)}[\Subst{s}{2}{x}]$.}{}

\iftag{defEx}{Pour déterminer si la !!{formula} existentielle
  $\lexists[x][!A(x)]$ est satisfaite, il faut vérifier si
  $\Sat{M}{\lnot \lforall[x][\lnot !A(x)]}[s]$. Ainsi,
  \[
  \Sat{M}{\lexists[x][(R(b,x) \lor R(x,b))]}[s].
  \]
  D'abord, $\Sat{M}{R(b,x) \lor R(x, b)}[\Subst{s}{1}{x}]$
  (la variante $\Subst{s}{3}{x}$ conviendrait aussi). Donc
  $\Sat/{M}{\lnot(R(b,x) \lor R(x,b))}[\Subst{s}{1}{x}]$, puis
  $\Sat/{M}{\lforall[x][\lnot(R(b,x) \lor R(x,b))]}[s]$ et,
  finalement, $\Sat{M}{\lnot\lforall[x][\lnot(R(b,x) \lor
  R(x,b))]}[s]$. En revanche,
  \[
  \Sat/{M}{\lexists[x][(R(b,x) \land R(x,b))]}[s].
  \]
  En effet, $\Sat{M}{\lforall[x][\lnot(R(b,x) \land
  R(x,b))]}[s]$, car il n'existe aucun $m \in \Domain M$ tel que
  $\Sat{M}{R(b,x) \land R(x,b)}[\Subst{s}{m}{x}]$. Ces exemples
  suggèrent que $\Sat{M}{\lexists[x][!A(x)]}[s]$ si et seulement si
  $\Sat{M}{!A(x)}[\Subst{s}{m}{x}]$ pour au moins un
  $m \in \Domain M$.}{}

\iftag{defAll}{Pour déterminer si la !!{formula} universelle
  $\lforall[x][!A(x)]$ est satisfaite, il faut vérifier si
  $\Sat{M}{\lnot\lexists[x][\lnot !A(x)]}[s]$. Par exemple,
  \[
  \Sat{M}{\lforall[x][(R(x,a) \lif R(a,x))]}[s],
  \]
  D'abord, $\Sat{M}{R(x,a) \lif R(a,x)}[\Subst{s}{m}{x}]$ pour tout
  $m \in \Domain M$ : $\Sat{M}{R(a,x)}[\Subst{s}{1}{x}]$, tandis que
  $\Sat/{M}{R(x,a)}[\Subst{s}{m}{x}]$ pour $m = 2$, $3$ ou~$4$.
  Il n'existe donc aucun $m \in \Domain M$ tel que
  $\Sat{M}{\lnot(R(x,a) \lif R(a,x))}[\Subst{s}{m}{x}]$;
  par suite, $\Sat/{M}{\lexists[x][\lnot(R(x,a) \lif R(a,x))]}[s]$.
  Nous obtenons ainsi
  $\Sat{M}{\lnot\lexists[x][\lnot(R(x,a) \lif R(a,x))]}[s]$.
  En revanche,
  \[
  \Sat/{M}{\lforall[x][(R(a,x) \lif R(x,a))]}[s],
  \]
  puisque $\Sat/{M}{R(a,x) \lif R(x,a)}[\Subst{s}{2}{x}]$ et donc
  $\Sat{M}{\lexists[x][\lnot(R(a,x) \lif R(x,a))]}[s]$. Ces exemples
  suggèrent que $\Sat{M}{\lforall[x][!A(x)]}[s]$ si et seulement si
  $\Sat{M}{!A(x)}[\Subst{s}{m}{x}]$ pour tout
  $m \in \Domain M$.}{}

Considérons maintenant un cas plus complexe :
\[
\lforall[x][(R(a,x) \lif \lexists[y][R(x,y)])].
\]
Puisque $\Sat/{M}{R(a,x)}[\Subst{s}{3}{x}]$ et
$\Sat/{M}{R(a,x)}[\Subst{s}{4}{x}]$, seuls les cas $m = 1$ et
$m = 2$ exigent d'examiner le conséquent de l'implication.
A-t-on $\Sat{M}{\lexists[y][R(x,y)]}[\Subst{s}{1}{x}]$ ? Il suffit
qu'il existe au moins un $n \in \Domain M$ tel que
$\Sat{M}{R(x,y)}[\Subst{\Subst{s}{1}{x}}{n}{y}]$. Pour $n = 1$,
nous avons $\Subst{\Subst{s}{1}{x}}{n}{y} = \Subst{s}{1}{y} = s$.
Comme $s(x) = 1$, $s(y) = 1$ et $\tuple{1,1} \in \Assign{R}{M}$,
la réponse est affirmative.

Reste à vérifier
$\Sat{M}{\lexists[y][R(x,y)]}[\Subst{s}{2}{x}]$.
Examinons les assignations de !!{variable}s
$\Subst{\Subst{s}{2}{x}}{n}{y}$. Pour $n = 1$, cette assignation
est~$s_2 = \Subst{s}{2}{x}$, qui ne satisfait pas $R(x,y)$ :
$s_2(x) = 2$, $s_2(y) = 1$ et
$\tuple{2,1}\notin \Assign{R}{M}$. En revanche, considérons
$\Subst{\Subst{s}{2}{x}}{3}{y} = \Subst{s_2}{3}{y}$.
Nous avons $\Sat{M}{R(x,y)}[\Subst{s_2}{3}{y}]$ puisque
$\tuple{2,3} \in \Assign{R}{M}$; donc
$\Sat{M}{\lexists[y][R(x,y)]}[s_2]$.

Ainsi, pour tout $m \in \Domain M$, ou bien
$\Sat/{M}{R(a,x)}[\Subst{s}{m}{x}]$ (si $m = 3$ ou~$4$), ou bien
$\Sat{M}{\lexists[y][R(x,y)]}[\Subst{s}{m}{x}]$ (si $m = 1$ ou~$2$).
Par conséquent,
\[
\Sat{M}{\lforall[x][(R(a,x) \lif \lexists[y][R(x,y)])]}[s].
\]
En revanche,
\[
\Sat/{M}{\lexists[x][(R(a,x) \land \lforall[y][R(x,y)])]}[s].
\]
La relation $\Sat{M}{R(a,x)}[\Subst{s}{m}{x}]$ ne vaut que pour
$m = 1$ ou $m = 2$. Pour chacun de ces deux cas, il existe pourtant
un $n \in \Domain M$ tel que
$\Sat/{M}{R(x,y)}[\Subst{\Subst{s}{m}{x}}{n}{y}]$ : on peut prendre
$n = 4$ lorsque $m = 1$, et $n = 1$ lorsque $m = 2$.
Ainsi $\Sat/{M}{\lforall[y][R(x,y)]}[\Subst{s}{m}{x}]$ pour
$m = 1$ et $m = 2$. Il n'existe donc aucun $m \in \Domain M$ tel que
$\Sat{M}{R(a,x) \land \lforall[y][R(x,y)]}[\Subst{s}{m}{x}]$.
\end{ex}

\iftag{defEx}{%
\begin{prop}\ollabel{prop:sat-ex}
$\Sat{M}{\lexists[x][!B(x)]}[s]$ si et seulement s'il existe une
variante en~$x$ de~$s$, notée~$s'$, telle que $\Sat{M}{!B(x)}[s']$.
\end{prop}

\begin{proof}
  Exercice.
\end{proof}
}{}

\tagprob{defEx}
\begin{prob}
  Démontrer la \olref[fol][syn][sat]{prop:sat-ex}.
\end{prob}
\tagendprob

\iftag{defAll}{%
\begin{prop}\ollabel{prop:sat-all}
$\Sat{M}{\lforall[x][!B(x)]}[s]$ si et seulement si, pour toute
variante en~$x$ de~$s$, notée~$s'$, $\Sat{M}{!B(x)}[s']$.
\end{prop}

\begin{proof}
  Exercice.
\end{proof}
}{}

\tagprob{defAll}
\begin{prob}
  Démontrer la \olref[fol][syn][sat]{prop:sat-all}.
\end{prob}
\tagendprob

\begin{prob}
Soit $\Lang L = \{c, f, A\}$, avec une !!{constant}, un
!!{function} unaire et un !!{predicate} binaire. Soit la
!!{structure}~$\Struct{M}$ définie par
\begin{enumerate}
\item $\Domain M = \{1, 2, 3\}$
\item $\Assign{c}{M} = 3$
\item $\Assign{f}{M}(1) = 2, \Assign{f}{M}(2) = 3, \Assign{f}{M}(3) = 2$
\item $\Assign{A}{M} = \{\tuple{1, 2}, \tuple{2, 3}, \tuple{3, 3}\}$
\end{enumerate}
(a) Soit $s(v) = 1$ pour toute !!{variable}~$v$. Déterminer si
\[
\Sat{M}{\lexists[x][(A(f(z), c) \lif \lforall[y][(A(y, x) \lor A(f(y),
      x))])]}[s]
\]
Justifier la réponse.

(b) Donner une autre !!{structure} et une assignation de
!!{variable}s dans lesquelles la !!{formula} n'est pas satisfaite.
\end{prob}

\end{document}
